\documentclass[tikz, border=10pt]{standalone}
\usepackage{pgfplots}
\pgfplotsset{compat=1.18}

\begin{document}

\begin{tikzpicture}
\begin{axis}[
    xlabel={Tensión $V_d$ [V]},
    ylabel={Corriente $I_d$ [A]},
    grid=major,
    width=12cm, height=8cm,
    xmin=0, xmax=0.8, 
    ymin=0, ymax=0.0004,
    scaled ticks=false, 
    /pgf/number format/1000 sep={},
    legend pos=north west,
    % Evita errores de memoria con la exponencial
    restrict y to domain=0:0.001 
]
    % 1. CURVA REAL (Simulación con N=1.984 e Is=14.11nA)
    \addplot[blue, thick] table[x index=0, y index=1, skip first n=1] {{TP_FE diodo.txt}}; 
    \addlegendentry{Curva Real (Simulada)}

    % 2. CURVA IDEAL (N=1)
    % Notarás que esta curva sube MUCHO antes que la azul
    \addplot[red, dashed, domain=0:0.6, samples=100] 
        {14.11e-9 * (exp(x / 0.02585) - 1)};
    \addlegendentry{Curva Ideal ($N=1$)}

    % 3. RECTA DE CARGA: (Vcc - Vd) / R
    \addplot[green!60!black, thick, domain=0:0.8] 
        {(3.5 - x) / 10150};
    \addlegendentry{Recta de Carga}

    % --- PUNTO Q ---
    \addplot[black, mark=*, mark size=2.5pt] coordinates {(0.5104, 0.0002945)};
    \node[anchor=south west, font=\small] at (axis cs:0.5104, 0.0002945) {\textbf{Q}};
    
\end{axis}
\end{tikzpicture}

\end{document}